\section{Einleitung}\label{einleitung}

In diesem Dokument möchte ich euch die Grundlagen von \LaTeX~erklären und wie
dieses Template genutzt werden kann. Dieses Dokument ist gleichzeitig mit
meinem Template erstellt und bietet einen kleinen Überblick, wie Hausarbeiten
formal ausgestaltet werden können. Aus diesem Grund hat diese Anleitung ein
formales Deckblatt und eine Eigenständigkeitserklärung. Um das Template zu
verstehen, schauen wir uns zu Beginn die Verzeichnisstruktur des
\LaTeX~Projektes an, um daraufhin den Aufbau der \texttt{hauptdatei.tex} zu
verstehen. Anschließend werden wir uns mit den verschiedenen Umgebungen wie
\texttt{figure}, \texttt{table} und \texttt{listings} beschäftigen. Dazu werden
wir uns auch mit Floatumgebungen in \LaTeX~auseinandersetzen müssen.
Abschließend blicken wir auf die Literaturverwaltung mit \texttt{BibTeX}. \\
Immer wieder werde ich dazu \LaTeX~Code zeigen und das Ergebnis dieses Codes.
Am Ende soll es euch möglich sein, mittels dieses Templates schöne und formal
korrekte Abgaben zu schreiben. Neben diesem Template existieren zusätzliche
Templates für die Erstellung von Postern, auf welche im Anhang eingegangen wird.\\
Grundsätzlich unterscheidet sich \LaTeX~von Word oder Pages dadurch, dass wir
Layout und Inhalt nicht zusammenführen. In \LaTeX~definieren wir unser Layout
über Befehle und Umgebungen und schreiben den Inhalt als \texttt{plain-text}.
So sehen wir beispielsweise, dass sich einige Wörter im obigen Text von anderen
unterscheiden. Dies geschieht in diesem Fall durch den Befehl
\verb|\texttt{beispieltext}|. Um den klassischen \LaTeX~Schriftzug zu erhalten,
kann man \verb|\LaTeX| nutzen. Es existieren eine Vielzahl solcher Kommandos, wir
werden uns nur mit den wenigsten und wichtigsten beschäftigen. Ziel dieser
Ausarbeitung soll es nicht sein, dass ihr zu absoluten \LaTeX~Expert:innen
werdet, sondern die Grundlagen beherrscht.\\ Meine Vorlagen sind so geschrieben,
dass sie ohne Probleme in der Informatikinfrastruktur gebaut werden können. 
Wollt ihr lokal mit diesen Vorlagen arbeiten, müsst ihr \LaTeX~installieren und
braucht zusätzlich noch einige Programme.  

